\chapter{Prototipo}\label{chap:prototipo}

El presente capítulo tiene como propósito describir el diseño, el alcance y la validación del prototipo funcional desarrollado. Para ello, se detallan los componentes que integran la primera versión, los datos empleados, los criterios de aceptación, la arquitectura implementada y el procedimiento definido para evaluar su desempeño. Finalmente, se presentan los escenarios de demostración y los resultados obtenidos en las pruebas iniciales.

\section{Alcance del prototipo funcional}
El prototipo funcional se orienta a validar el núcleo técnico de la propuesta. Para ello, contempla un sistema capaz de procesar estados de cuenta bancarios mediante una API REST desarrollada en FastAPI, la cual implementa un flujo automatizado de extracción basado en un extractor híbrido con OCR y un modelo de lenguaje, diseñado para identificar los campos clave requeridos. Asimismo, incorpora una base de datos SQLite para el registro de extracciones, correcciones manuales y métricas de desempeño. En conjunto, estos componentes permiten materializar la solución a nivel técnico, asegurar la trazabilidad de los resultados y generar evidencia sobre el comportamiento del sistema durante su validación.


\section{Construcción del prototipo}

La construcción del prototipo se basa en una arquitectura multinivel que organiza el sistema en capas con responsabilidades diferenciadas. Este enfoque facilita la separación de funciones, el mantenimiento del código y la trazabilidad del proceso, lo que resulta fundamental para asegurar que cada componente del sistema pueda evolucionar de manera ordenada sin comprometer el funcionamiento integral de la solución.

La capa de presentación se implementa con \textit{FastAPI} \footnote{Documentación oficial de FastAPI: \url{https://fastapi.tiangolo.com/}} y utiliza \textit{Pydantic} \footnote{Documentación oficial de Pydantic: \url{https://docs.pydantic.dev/latest/}} para la validación estructural de los datos. Además, expone los \textit{endpoints} que permiten cargar archivos PDF, registrar correcciones, consultar los registros del sistema y obtener métricas de ejecución. Su función principal es recibir las solicitudes externas y estandarizar los datos de entrada y salida. Por su parte, la capa de dominio concentra la lógica principal del prototipo, ya que contiene los servicios de extracción, las reglas necesarias para el procesamiento de los archivos y los modelos que representan las entidades involucradas en el proceso. Esta capa opera de manera independiente de la infraestructura técnica, lo que refuerza la modularidad del sistema y favorece la reutilización de la lógica de negocio.

La capa de infraestructura proporciona los mecanismos necesarios para ejecutar las operaciones del sistema. Incluye la interacción con la base de datos SQLite, la integración con el OCR y la gestión de los archivos temporales generados durante el procesamiento. A su vez, la capa de aplicación centraliza la configuración general del sistema, administra el manejo de excepciones y controla el sistema de \textit{logging}, coordinando así la inicialización y operación de los distintos componentes del prototipo. En este contexto, el extractor principal, denominado \textit{ClaudeOCRExtractor}, integra un motor OCR tradicional, Tesseract, con las capacidades de análisis visual y extracción estructurada del modelo \textit{Claude AI}. Esta combinación permite procesar documentos con variaciones de formato y obtener resultados más consistentes, mientras que las validaciones automáticas posteriores, como la verificación de la longitud de la CLABE, la normalización del nombre del banco y la revisión de consistencia entre campos, fortalecen la confiabilidad de la información extraída.

Finalmente, para asegurar la trazabilidad de los resultados y evaluar de manera objetiva el desempeño del sistema, el prototipo incorpora una tabla de registros que almacena, para cada documento procesado, los valores extraídos, las correcciones manuales asociadas, así como los aciertos y errores identificados. A partir de esta información, se calculan métricas automáticas, entre ellas la precisión por campo, la proporción de correcciones y las variaciones de desempeño a lo largo del tiempo. Esto permite no solo monitorear el comportamiento del prototipo, sino también generar evidencia objetiva sobre su desempeño y facilitar la identificación de oportunidades de mejora continua.

\section{Verificación preliminar del prototipo}

Como verificación inicial del funcionamiento del prototipo, se realizaron pruebas sobre un conjunto de cinco archivos PDF. En esta demostración, el prototipo procesó correctamente el 100\% de los documentos sin errores críticos. Del total evaluado, el 60\% obtuvo una extracción completa de los tres campos analizados, mientras que el 40\% restante presentó una extracción parcial, con al menos un campo correctamente identificado. La precisión por campo evidenció un desempeño particularmente sólido en la identificación del banco, correcta en todos los casos. Por su parte, tanto la CLABE como la identificación del titular alcanzaron un 60\% de aciertos. En conjunto, el 40\% de los campos extraídos requirió corrección manual, mientras que el tiempo promedio de procesamiento se situó entre 15 y 25 segundos por archivo. Estos resultados muestran que el prototipo opera de manera estable sobre documentos reales y ofrece evidencia preliminar de su factibilidad técnica. Cabe señalar que este extractor basado en OCR fue posteriormente superado por el enfoque de envío directo de PDF al modelo, como se detalla en el capítulo \ref{chap:experimentos}.
